\section{The Five-Vector Taxonomy}
\label{sec:taxonomy}

We define a \emph{gaming vector} as any degree of freedom available to
the actor responsible for deploying the redistricting algorithm that can
be varied to influence the partisan composition of the output.
Gaming vectors are distinct from design choices: a design choice is a
single, publicly documented decision made before any partisan outcome is
known; a gaming vector requires the hostile actor to observe partisan
implications before exercising the freedom.

\begin{definition}[Gaming Vector]
A \emph{gaming vector} is a tuple $(X, \Delta D, \delta)$ where:
$X$ is the degree of freedom that can be manipulated;
$\Delta D$ is the maximum achievable partisan shift (measured in expected
Democratic seats nationally) when $X$ is optimized for partisan advantage;
and $\delta$ is the detectability---the probability that the manipulation
is detected under the \dia{} audit protocol.
\end{definition}

Table~\ref{tab:taxonomy} summarises the five vectors, their attack surfaces,
estimated maximum effects, and detection mechanisms.

\begin{table}[h]
\centering
\caption{The five gaming vectors. $\Delta D$: maximum expected partisan
         shift, national 435-seat delegation. Detectability: probability of
         detection under the \dia{} audit protocol after pre-registration.
         Vectors V1--V3 affect the plan geometry; V4--V5 affect only presentation.}
\label{tab:taxonomy}
\begin{tabular}{clccc}
\toprule
ID & Gaming Vector & $\Delta D$ & Detectability & Plan-affecting? \\
\midrule
V1 & Parameter manipulation & $\leq 0.3$ & 1.00 & Yes \\
V2 & Input data manipulation & $\leq 0.4$ & $\approx 1.00$ & Yes \\
V3 & Geographic definition gaming & $\leq 1.5$ & 0.95 & Yes \\
V4 & Selective result presentation & N/A & 0.80 & No \\
V5 & Algorithm version selection & $\leq 0.3^*$ & 1.00 & Yes (indirectly) \\
\midrule
\multicolumn{2}{l}{V1--V3 combined (additive upper bound)} & $\leq 2.2$ & & \\
\multicolumn{2}{l}{V1--V3 combined (empirical)} & $\leq 0.5$ & & \\
\bottomrule
\end{tabular}
\begin{flushleft}
\small $^*$Within the same algorithm family (GeoSection); cross-family
differences are $\leq 18$ seats but are constrained by B.0's structure pre-specification.
\end{flushleft}
\end{table}

\subsection{Classification by Plan Impact}

Vectors V1, V2, and V5 affect the plan geometry because they change the
computational process that generates the plan.
Vector V3 (geographic definition gaming) affects the plan geometry through
the adjacency graph structure.
Vector V4 (selective result presentation) does not change the plan geometry
at all; it only changes which plan is reported to adjudicators.

The distinction is important for \dia{} design.
V4 and V5 are controlled entirely by disclosure requirements and binary
pinning; they do not require any analysis of the computational outputs.
V1--V3 require that the computation be reproducible from independently
verifiable inputs.

\subsection{Interaction Between Vectors}

A sophisticated adversary might ask: can multiple vectors be combined
to amplify the partisan effect beyond the individual maxima?
Section~\ref{sec:combined} shows that the answer is no---the interaction
structure of the \bisect{} algorithm ensures that the combined effect is
subadditive.
The geometric reason is that parameter effects operate on leaf-level
METIS cuts within a fixed bisection tree structure; input data effects
change the vertex weights within the same graph; geographic effects change
the graph topology.
All three effects converge on the same geometric bottleneck---the
competitive-district geography---and cannot compound each other.
